% ================================================
% 文件: 08_软件定义无线电.tex
% 章节: 第15章 软件定义无线电
% 所属部分: 接收机技术
% 版本: v1.0
% ================================================
% 本章目录结构（树状）:
% \chapter{软件定义无线电}
% ├── \section{软件定义无线电概述}
% ├── \section{SDR基本架构}
% ├── \section{SDR关键技术}
% ├── \section{SDR接收机设计要点}
% └── \section{SDR应用与发展趋势}
% ================================================

\chapter{软件定义无线电}
\label{chap:sdr}

\section{软件定义无线电概述}
\label{sec:sdr_intro}

软件定义无线电（Software Defined Radio，SDR）是一种其关键功能尽可能由软件实现的无线电通信系统。它的核心思想是：把传统无线电中由固定硬件电路完成的调制与解调、滤波、频率变换、编码与解码等功能，转移到数字域中用软件算法完成；模数转换器（ADC）与数模转换器（DAC）的位置尽可能靠近天线，使射频信号尽早数字化，从而让同一套通用硬件平台通过加载不同的软件（波形）来扮演完全不同的电台。1992 年，Joseph Mitola 在 IEEE NEC 会议与相关论文中系统提出"软件无线电"的概念与架构，奠定了这一领域的理论基础；此后"软件定义无线电"作为更工程化的称谓被广泛采用。

SDR 的发展与数字器件的进步紧密相连。20 世纪 70 年代，美国国防高级研究计划局（DARPA）出于多军种电台互联互通的需要开始探索通用无线电平台，早期研究集中在军用领域；80 年代，数字信号处理器（DSP）的成熟催生了第一批以 DSP 为核心的软件无线电原型，验证了用软件完成解调与滤波的可行性。进入 90 年代，高速 ADC/DAC 与现场可编程门阵列（FPGA）的实用化，使宽带采样与实时数字前端成为现实，SDR 由概念走向工程产品。2000 年以后，SDR 进入标准化与商业化阶段，ITU 等组织推动了相关标准研究；与此同时，以 GNU Radio 为代表的开源软件生态和以通用软件无线电外设为代表的低成本硬件，把 SDR 从昂贵的军用设备变为科研、教学与业余无线电领域人人可用的工具。近年来，随着 5G/6G、物联网与人工智能技术的发展，SDR 已从专用平台演进为通信系统研发的标准基础设施。

与传统硬件无线电相比，SDR 的根本差异在于"功能由谁决定"：前者一旦硬件电路定型，功能即告固定，支持新制式必须重新设计硬件；后者的射频前端保持相对通用，功能随软件而变，同一台设备今天可以是广播接收机，明天可以是频谱监测仪。两者的典型对比如下表所示。

\begin{table}[htbp]
\centering
\caption{传统硬件无线电与软件定义无线电的对比}
\label{tab:sdr_compare}
\begin{tabular}{|p{2.2cm}|p{5.2cm}|p{6.2cm}|}
\hline
\textbf{对比项} & \textbf{传统硬件无线电} & \textbf{软件定义无线电} \\
\hline
功能实现 & 调制解调、滤波等由专用模拟/数字硬件电路固定实现 & 同类功能由软件算法（FPGA/DSP/GPP）实现 \\
\hline
支持新制式 & 需重新设计并更换硬件，周期长、成本高 & 加载或升级波形软件即可，周期短 \\
\hline
硬件复用 & 一种硬件通常对应一种用途 & 同一平台可配置为接收机、发射机、监测仪等 \\
\hline
升级方式 & 更换模块或整机 & 更新软件版本，必要时重配 FPGA \\
\hline
成本结构 & 多制式需多套设备，重复投入 & 平台通用化，初期较高、总体摊薄 \\
\hline
典型短板 & 灵活性差，难以适应标准演进 & 对 ADC/DAC、时钟与处理能力要求高 \\
\hline
\end{tabular}
\end{table}

概括而言，SDR 具有以下显著特点：一是\textbf{灵活性}，通过修改软件即可适应不同的通信标准与协议，无需更换硬件；二是\textbf{可重构性}，同一设备可在不同时刻承担不同任务，如接收广播、进行业余无线电通信或实施频谱监测；三是\textbf{易于升级}，当出现新的通信标准或需要改进性能时，升级软件即可完成；四是\textbf{高度集成}，多种通信功能集中于一个平台，减少了设备数量与系统复杂性；五是\textbf{成本效益}，通用平台摊薄了多制式应用的成本，虽然初始投入可能较高，但长期维护与升级成本显著降低。这些特点使 SDR 成为现代通信设备，特别是多制式、多任务设备的首选架构。

\section{SDR基本架构}
\label{sec:sdr_architecture}

SDR 系统由天线、射频前端、ADC/DAC 与数字信号处理平台四大部分组成，其设计目标是把模拟部分压缩到"不得已而为"的最小范围。以接收方向为例，信号链路为：天线接收射频信号后，先经预选滤波器抑制带外强干扰，由低噪声放大器（LNA）放大，再经过频率变换（或直接采样）送入 ADC；ADC 输出的高速数字序列首先进入数字下变频（DDC）环节，由数控振荡器（NCO）产生的本振序列与信号相乘，将目标信号搬移到数字基带，随后经过抽取滤波（CIC、半带、FIR 级联）降低数据速率；最后由基带处理软件完成信道滤波、同步、解调与解码。发射链路与接收链路对称：基带软件生成调制信号，经数字上变频（DUC）提高采样率后送入 DAC，再经模拟上变频与功率放大由天线辐射。

\textbf{采样与奈奎斯特准则}。ADC 是 SDR 的"咽喉"，采样定理决定了它的工作边界：对最高频率为 $f_{\max}$ 的低通信号，采样率必须满足
\[
f_s \ge 2f_{\max},
\]
才能无失真地恢复原信号。直接对射频信号采样需要极高的采样率与模拟输入带宽，因此在多数工程实现中，仍保留一级模拟混频把信号搬移到合适的中频或基带，再进行数字化。

\textbf{带通采样（欠采样）}。对带限于 $[f_L, f_H]$ 的带通信号（带宽 $B = f_H - f_L$），并不要求采样率高于 $2f_H$：只要采样率落入允许区间，信号落在某一奈奎斯特区内，混叠就有害变有用——采样过程本身完成了下变频。允许的采样率区间可写为
\[
\frac{2f_H}{m} \le f_s \le \frac{2f_L}{m-1}, \qquad m = 1, 2, \ldots, \left\lfloor \frac{f_H}{B} \right\rfloor,
\]
其中 $m$ 为采样后信号所在的奈奎斯特区编号。带通采样使 ADC 有可能直接对高中频甚至射频信号采样，是"射频直接采样"架构的理论基础，但区间上、下限随 $m$ 增大而收窄，对采样时钟的准确度与模拟带通滤波器的矩形度要求苛刻。

\textbf{过采样与处理增益}。以远高于奈奎斯特率的速率采样，可以把量化噪声功率摊到更宽的数字频带内，再经数字滤波与抽取去除带外部分。设信号带宽为 $B$，采样率为 $f_s$，则相对两倍带宽的临界采样，可获得的处理增益为
\[
G_p = 10\lg\frac{f_s}{2B} \ \mathrm{dB}.
\]
过采样每提高一倍，信噪比约改善 3 dB，相当于增加约半个有效位。这一原理也是 $\Sigma$-$\Delta$ 型 ADC 以低速高位数结构实现高精度采样的基础。

\textbf{数字下变频}。DDC 是数字前端的核心。设 ADC 输出序列为 $x(n)$，需要接收的载波频率对应数字角频率 $\omega_0$，则数字混频可表示为
\[
y(n) = x(n)\, e^{-j\omega_0 n},
\]
其中复指数序列由 NCO 查表产生。混频后的信号经低通滤波与 $R$ 倍抽取，数据速率降为 $f_s/R$，供基带处理。抽取滤波通常采用"CIC 粗调 + 半带滤波 + FIR 精调"的级联结构，以在降低运算量的同时保证带外抑制。快速傅里叶变换（FFT）的频率分辨率为
\[
\Delta f = \frac{f_s}{N_{\mathrm{FFT}}},
\]
它决定了频谱显示与信号搜索的最小分辨能力。

SDR 接收链路各环节的功能与关键指标汇总如下表。

\begin{table}[htbp]
\centering
\caption{SDR 接收链路主要环节}
\label{tab:sdr_chain}
\begin{tabular}{|p{2.6cm}|p{5.4cm}|p{5.6cm}|}
\hline
\textbf{环节} & \textbf{主要功能} & \textbf{关键指标} \\
\hline
预选滤波 & 抑制带外强干扰，防止阻塞与混叠 & 插入损耗、带外抑制、可调谐性 \\
\hline
低噪声放大 & 在加入尽量少噪声的前提下放大微弱信号 & 噪声系数、增益、线性度（IP3） \\
\hline
频率变换 & 将射频搬移到中频或基带（或省略，直接采样） & 变频损耗、镜像抑制、本振相噪 \\
\hline
抗混叠滤波 & 限制 ADC 输入带宽，保证采样不产生带内混叠 & 阻带衰减、过渡带宽度 \\
\hline
ADC & 模拟信号数字化 & 采样率、位数、ENOB、SFDR \\
\hline
数字下变频 & NCO 混频、滤波、抽取，降低数据速率 & 抽取比、带外抑制、相位噪声 \\
\hline
基带处理 & 同步、解调、解码与协议处理 & 处理时延、吞吐率、可重构性 \\
\hline
\end{tabular}
\end{table}

\section{SDR关键技术}
\label{sec:sdr_key_technologies}

\textbf{宽带射频前端}。SDR 要"一套前端走天下"，射频前端必须在宽频段内同时满足匹配、噪声与线性度三方面要求。宽带阻抗匹配意味着天线与放大器要在数个倍频程内保持良好驻波；宽带线性度意味着在多强信号同时进入时，互调产物仍须足够低；宽带噪声系数则要求放大器在宽频段内保持低噪。为此，SDR 前端普遍采用可调谐（或开关切换的子倍频程）预选滤波器组压缩瞬时带宽，配合宽带 LNA 与高线性度混频器，并利用 AGC 与增益分配策略防止大信号使后级饱和。

\textbf{高速 ADC/DAC}。转换器决定了 SDR 性能的上限。理想 $N$ 位转换器的量化信噪比为
\[
\mathrm{SNR}_{\mathrm{ideal}} = 6.02N + 1.76 \ \mathrm{dB},
\]
即每增加 1 位，动态范围约提高 6 dB。实际器件受噪声与失真影响，有效位数（ENOB）低于标称位数：
\[
\mathrm{ENOB} = \frac{\mathrm{SNR}_{\mathrm{dB}} - 1.76}{6.02}.
\]
除量化噪声外，采样时钟的孔径抖动 $t_j$ 对高频输入信号的 SNR 构成硬限制：
\[
\mathrm{SNR}_{\mathrm{jitter}} = -20\lg\left(2\pi f_{\mathrm{in}} t_j\right) \ \mathrm{dB},
\]
输入频率越高，抖动的影响越显著，这也是高速 SDR 对时钟质量极为敏感的原因。衡量转换器大信号能力的另一指标是无杂散动态范围（SFDR），它决定了强信号旁同时接收弱信号的能力。

\textbf{数字处理平台}。SDR 的处理负载通常分层分配：FPGA 擅长并行、确定性时延的流式运算，承担数字上下变频、抽取滤波、相关与交织等高吞吐任务；DSP/GPP 承担调制解调、信道编码等逐符号或逐帧处理；GPU 与通用 CPU 的 SIMD 指令则在离线分析与软件化基站中发挥作用。三种平台的取舍在于吞吐率、功耗与开发周期的平衡，工程上常见"FPGA 打前端、CPU 打上层"的组合。

\begin{table}[htbp]
\centering
\caption{SDR 数字处理平台对比}
\label{tab:sdr_platforms}
\begin{tabular}{|p{2.4cm}|p{4.6cm}|p{4.2cm}|p{2.6cm}|}
\hline
\textbf{平台} & \textbf{特点} & \textbf{适用任务} & \textbf{相对功耗} \\
\hline
FPGA & 并行流水、时延确定、可重构硬件逻辑 & 数字前端、高速滤波、相关捕获 & 中 \\
\hline
DSP & 乘加结构高效、实时性好 & 调制解调、信道编译码 & 中 \\
\hline
GPP/GPU & 开发生态好、灵活性最高 & 协议栈、频谱分析、离线处理 & 高 \\
\hline
\end{tabular}
\end{table}

\textbf{软件架构与生态}。典型的 SDR 软件栈分为三层：底层为设备驱动与抽象层（如面向通用软件无线电外设的 UHD 驱动框架），向上屏蔽具体硬件差异；中间层为信号处理框架，以 GNU Radio 为代表，把处理单元抽象为"块"（Block），用"流图"（Flowgraph）描述数据流，开发者通过连接现成模块或编写新模块快速搭建收发系统；应用层则是具体业务，如广播接收、数字模式通信、信号分析与协议实现。此外，MATLAB/Simulink 常用于算法建模与验证，各类专用软件（如数字模式通信与卫星遥测软件）则面向最终用户。波形加载与重配置机制是 SDR 的灵魂：设备上电后加载固件与 FPGA 镜像，再按任务装载波形描述文件，即可在"多种电台"之间切换。

\textbf{时钟与同步}。采样时钟与本振的相位噪声直接决定接收机的带内噪声本底与多普勒容限；频率准确度与稳定度影响免搜索接收与多通道一致性。对频率准确度要求高的应用常采用 GPS 驯服恒温晶振（GPSDO）作参考；多通道系统还须保证采样时钟同源与触发同步，否则波束形成、干涉测向与多站协同无法实现。

\textbf{非理想校正}。零中频与直接变频架构会引入直流偏置、I/Q 幅相失衡与本振泄漏，数字域可通过直流跟踪与自适应失衡校正算法抑制；时钟泄漏与电源纹波则需在板级设计中配合去耦与屏蔽解决。这些"数字补模拟"的手段是 SDR 得以用通用器件逼近专用接收机性能的关键。

\section{SDR接收机设计要点}
\label{sec:sdr_receiver_design}

\textbf{架构选择}。SDR 接收机在三种主流架构中取舍：其一为\textbf{射频直接采样}，信号经滤波放大后直接进入 ADC，链路最短、灵活性最高，但受 ADC 模拟带宽、采样率与抖动限制，目前多用于中短波、VHF/UHF 等较低频段；其二为\textbf{零中频（直接变频）}，射频一步变频到基带后采样，省去中频滤波与镜频抑制，集成度最高，但须处理直流偏置与闪烁噪声；其三为\textbf{超外差加数字中频}，模拟域完成一次变频与主要滤波，在中频数字化后交由数字前端处理，性能最稳妥，是高性能监测与通信接收机的常见选择。三种架构的差别本质上仍是第 8 章所述超外差原理在数字时代的延伸：把"哪一级用模拟、哪一级用数字"的边界推移到哪里。

\textbf{混叠与镜像抑制}。直接采样架构中，凡落在采样带宽之外或过渡带内的信号都会按
\[
f_{\mathrm{alias}} = \left| f_{\mathrm{sig}} - k f_s \right|
\]
折叠回带内形成混叠干扰，因此抗混叠滤波器的阻带抑制必须覆盖最强的带外信号电平；带通采样架构还须防止其他奈奎斯特区的强信号"串区"进入目标区。超外差加数字中频架构中，镜像频率抑制由模拟前端承担，数字域则需抑制本振泄漏与半像频干扰。

\textbf{噪声与灵敏度}。SDR 接收机的灵敏度仍由级联噪声系数决定：
\[
F = F_1 + \frac{F_2 - 1}{G_1} + \frac{F_3 - 1}{G_1 G_2} + \cdots,
\]
设计时应保证第一级 LNA 具有足够增益，使 ADC 与其后级的噪声贡献被压低。ADC 本身可折算为一级"等效噪声源"：把量化噪声与热噪声合并后换算成等效噪声系数参与级联计算。与第 10 章的测试方法一致，SDR 接收机的灵敏度仍以最小可检测信号（MDS）或信纳比（SINAD）表征，但须注意数字处理增益（如匹配滤波、长积分）可以换取灵敏度，而 AGC 与数字增益分配不当则会侵蚀它。

\textbf{动态范围设计}。宽带采样使"多信号共存"成为常态，动态范围成为 SDR 最受考验的指标：一是\textbf{阻塞与 AGC}，强带外信号经前端压缩后不得淹没带内弱信号，需在模拟域尽早衰减；二是\textbf{互调}，双音三阶互调产物由前端 IP3 与 ADC 的 SFDR 共同决定，增益分配应使各级等效 IP3 与噪声贡献达到平衡；三是\textbf{数字域的余量}，抽取滤波的溢出保护、浮点或高位宽定点运算的选择，都会影响弱信号最终可提取的程度。

\textbf{选型与验证}。选择 SDR 平台时应关注：频率覆盖与瞬时带宽是否满足任务；ADC/DAC 的位数、采样率与 ENOB、SFDR；时钟方案（内部晶振质量、外部参考输入、GPSDO 支持）；接口与主机吞吐（USB/GbE/10GbE/PCIe）能否承载实时采样流；软件生态与开源支持程度；以及收发能力与功率等级是否匹配用途。装机后应按第 10 章的方法复测灵敏度、选择性、动态范围与频率准确度，并与理论估算对照，确认前端增益分配与时钟质量达到设计预期。

\section{SDR应用与发展趋势}
\label{sec:sdr_applications}

SDR 的应用几乎覆盖无线电的全部场景。\textbf{在广播接收方面}，同一台 SDR 配合不同软件即可接收 AM、FM 及数字广播，并可对信号进行录音与离线分析，成为广播收听与信号研究的利器。\textbf{在业余无线电方面}，SDR 支持短波数字模式通信、卫星收发与远程电台（网络化遥控电台）建设，其宽频谱显示能力也改变了传统的"旋钮调谐"操作方式。\textbf{在频谱监测方面}，宽带 SDR 可实现全景扫描、信号截获与测向配合，为频谱管理与干扰查处提供低成本手段，其原理与第 12 章监测接收机一脉相承。\textbf{在通信研发方面}，SDR 是 4G/5G 及未来 6G 空口技术快速原型验证的标准平台，配合开源基站软件可实现协议栈级的实验。\textbf{在教学与科研方面}，低成本 SDR 让学生能够亲手"打开"真实的无线电信号，把抽象的采样定理、调制理论与滤波设计变为可视的实验。

常见入门与专业平台各具定位：RTL-SDR 基于 DVB-T 接收芯片被改造为宽带接收器，价格极低、适合入门与广播监测，但仅能接收且位数与带宽有限；HackRF One 为开源半双工收发平台，频段覆盖宽，适合协议实验与逆向研究；USRP 系列面向科研与产业原型，采样率高、生态完善；LimeSDR、BladeRF 等则在带宽与价格之间提供多种选择。下表给出常见平台的典型定位（具体参数以厂商最新资料为准）。

\begin{table}[htbp]
\centering
\caption{常见 SDR 平台及其典型定位}
\label{tab:sdr_devices}
\begin{tabular}{|p{2.6cm}|p{2.4cm}|p{3.4cm}|p{5.2cm}|}
\hline
\textbf{平台} & \textbf{收发能力} & \textbf{典型频段} & \textbf{特点与典型用途} \\
\hline
RTL-SDR & 仅接收 & 约 24MHz$\sim$1.7GHz & 成本极低、8 位，入门学习、广播与频谱观察 \\
\hline
HackRF One & 半双工收发 & 约 1MHz$\sim$6GHz & 开源硬件、8 位，协议实验与教学演示 \\
\hline
USRP 系列 & 收发 & 约直流$\sim$6GHz & 高采样率、生态完善，科研与通信系统原型 \\
\hline
LimeSDR / BladeRF & 收发 & 约直流$\sim$6GHz & 宽带多通道，介于入门与专业之间的选择 \\
\hline
\end{tabular}
\end{table}

展望未来，SDR 的发展呈现四个方向：一是\textbf{更靠前的数字化}，随着 ADC 性能提升与数字前端技术进步，直接射频采样的频段不断上移，模拟链路进一步缩短；二是\textbf{智能化}，机器学习被用于调制识别、信号分选与频谱预测，SDR 开放的数据接口为此提供了天然试验床，也推动认知无线电与动态频谱共享走向实用；三是\textbf{网络化与开放化}，软件化基站与开放无线接入网（O-RAN）把 SDR 思想推广到移动通信基础设施，通用硬件加开放软件成为产业趋势；四是\textbf{安全与可靠性}，波形与固件的可更新性在带来灵活性的同时也引入供应链与信息安全问题，安全启动、签名升级与射频域抗干扰设计正成为 SDR 工程的新课题。

总体而言，SDR 并未否定超外差等经典原理，而是把它们装入通用平台并以软件重新表达。掌握模拟射频与数字信号处理两条主线，理解"边界放在哪里"的工程权衡，是用好与设计 SDR 的关键。
